ubsection{Karl Pearson}
		
			O Coeficiente de Correlação de Karl Pearson (1857-1936), Coeficiente de Correlação Produto-Momento ou Densidade $ \rho $ (letra grega) de Pearson, obviamente retorna o grau de correlação entre dois conjuntos de dados ($X$, $Y$). Em outras palavras, o coeficiente retorna a direção da tendência (crescente ou decrescente) de um conjunto em relação ao outro e vice-versa.
			
			Este coeficiente pode assumir valores entre $-1$ e $1,$ inclusive algum dos extremos \cite{pearson1895}.
			
			\begin{align}
				\rho & = \dfrac{cov(X, Y)}{\sqrt{var(X)var(Y)}}
			\end{align}
			
			Sendo: \\
			
			$ \begin{cases}
					\rho & \text{: Coeficiente de Correlação de Karl Pearson} \\[0.3cm]
					cov(X, Y) & \text{: covariância entre os conjuntos $X$ e $Y$} \\[0.3cm]
					var(X) & \text{: variância do conjunto $X$} \\[0.3cm]
					var(Y) & \text{: variância do conjunto $Y$}
				\end{cases} $
		
